% Part: first-order-logic
% Chapter: beyond
% Section: introduction

\documentclass[../../../include/open-logic-section]{subfiles}

\begin{document}

\olfileid{fol}{byd}{int}

\olsection{అవలోకనం}

మనకు ఆసక్తికరమైన ఏకైక తర్కవ్యవస్థ మొదటిస్థాయి తర్కం కాదు:
మొదటిస్థాయి తర్కానికి అనేక విస్తరణలు, భేదాలు ఉన్నాయి. ఒక తర్కం
సాధారణంగా ఒక భాషకు ఔపచారిక నిర్దేశం, సాధారణంగా---కానీ ఎల్లప్పుడూ
కాదు---ఒక నిగమన వ్యవస్థ, సాధారణంగా---కానీ ఎల్లప్పుడూ కాదు---ఒక
ఉద్దేశిత అర్థవిచారం కలిగి ఉంటుంది. అయితే ఈ పదాన్ని సాంకేతికంగా
వాడటం ఒక స్పష్టమైన ప్రశ్నను లేవనెత్తుతుంది: అంతఃప్రజ్ఞాత్మక లేదా
తాత్త్విక అర్థంలో వాడే ``తర్కం'' అనే మాటతో, మొదటిస్థాయి తర్కం కాని
తర్కాలకు ఏమి సంబంధం? కింద వివరించే వ్యవస్థలన్నీ ఏదో ఒక రకమైన
హేతుచింతనను నమూనీకరించడానికి రూపొందించినవే; వాటిని తార్కికమైనవిగా
చేసేది ఏమిటో చెప్పగలమా?

దీనికి సులభమైన సమాధానాలు లేవు. ``తర్కం'' అనే మాటను వేర్వేరు
సందర్భాల్లో వేర్వేరు విధాలుగా వాడుతారు; ``సత్యం'' అనే భావనలాగే ఈ
భావననూ అనేక తాత్త్విక దృక్కోణాల నుంచి విశ్లేషించారు. ఉదాహరణకు, ఏ
ప్రకటనలు అవశ్యంగా సత్యమో, అనుభవానికి ముందే సత్యమో, అతార్కిక పదాల
అర్థనిర్దేశంతో సంబంధం లేకుండా సత్యమో, వాటి రూపం వల్ల సత్యమో, లేదా
భాషా సంప్రదాయం వల్ల సత్యమో నిర్ణయించడమే తార్కిక హేతుచింతన లక్ష్యమని
భావించవచ్చు; ఈ భావనల్లో ప్రతిదానికీ ఎంతో వివరణ అవసరం. గణితంలో వాడే
తర్కంపైనే దృష్టిని పరిమితం చేసినా, దాని పరిధి విషయంలో ఏకాభిప్రాయం
చాలా తక్కువ. ఉదాహరణకు, ఫ్రేగె ప్రారంభించిన {\em తర్కవాద} సంప్రదాయంలో
రసెల్, వైట్‌హెడ్‌లు \textit{Principia Mathematica}లో తర్కాన్ని
పునాదిగా చేసుకొని గణితాన్ని అభివృద్ధి చేయడానికి ప్రయత్నించారు. వారి
తర్కవ్యవస్థ కింద వివరించే వ్యవస్థను పోలిన ఉన్నత-టైపు తర్క రూపం.
చివరికి, చాలా ప్రమాణాల ప్రకారం పూర్తిగా తార్కికమైనవిగా అనిపించని
స్వీకృతాలను (ముఖ్యంగా అనంతతా స్వీకృతాన్ని, అవకరణీయతా స్వీకృతాన్ని)
వారు ప్రవేశపెట్టవలసి వచ్చింది. అయినప్పటికీ ఉన్నత-స్థాయి హేతుచింతనలో
కొన్ని రూపాలను తార్కికమైనవిగా అంగీకరించాలని భావించవచ్చు. దీనికి
భిన్నంగా, తన సత్తావిచారంలో ``ప్రతిజ్ఞావాక్యాలను'' చర్చకు తగిన
వస్తువులుగా అంగీకరించని క్వైన్, ద్వితీయ-స్థాయి, ఉన్నత-స్థాయి తర్కాలు
నిజానికి గొర్రెతోలు కప్పుకున్న సమితి సిద్ధాంత రూపాలేనని వాదిస్తాడు;
మరో మాటలో, విధేయాలపై పరిమాణీకరణ ఉన్న వ్యవస్థలు పూర్తిగా తార్కికమైనవి
కావు.

ప్రస్తుతానికి ఇటువంటి తాత్త్విక సమస్యలను మరొక సందర్భానికి వదిలి,
కిందివి తార్కికమైనవైనా కాకపోయినా, వివిధ రకాల హేతుచింతనలకు ఔపచారిక
ఆదర్శీకరణలని భావించడం మేలు.

\end{document}
